\section{Kierunki dalszych badań}

Najbardziej naturalnym kierunkiem jest rozwój modeli hybrydowych. Wyniki fazy~II wskazują, że największy zysk daje planowanie sekwencji akcji, a rozszerzenie do 28 cech daje drugi, co do siły, efekt wśród modyfikacji. Szczególnie obiecujący wydaje się wariant \textit{28 głębokościowy} --- połączenie przeszukiwania głębokościowego z rozszerzoną funkcją oceny. W tej pracy nie zbadano tej kombinacji, ale logicznie łączy ona oba najsilniejsze źródła poprawy jakości.

Drugi kierunek to poszerzenie warunków oceny turniejowej: większy zbiór talii oraz konfrontacja z większym gronem agentów zewnętrznych, jak np. MCTS lub inne warianty lookahead. Takie testy pozwoliłyby ocenić odporność modyfikacji na zmianę środowiska eksperymentu i ograniczyć ryzyko dopasowania.

Trzeci kierunek to pogłębienie analizy procesu uczenia --- przejście od podsumowań turniejów do analizy decyzji w konkretnych turach, identyfikacji kluczowych stanów planszy i śledzenia wpływu grup wag na wybór akcji.

Czwarty kierunek to dalsza optymalizacja algorytmiczna. Oprócz rodziny SHADE warto przetestować inne metody adaptacyjne i podejścia wielokryterialne, w których funkcja celu uwzględnia nie tylko skuteczność turniejową, ale też stabilność zachowania agenta oraz koszt obliczeniowy podejmowania decyzji. Osobno warto zbadać \textit{linkage learning} --- uczenie powiązań między cechami funkcji oceny, które mogłoby ułatwić optymalizację w przestrzeniach o wyższym wymiarze, takich jak reprezentacja 63-parametrowa.